\documentclass[conference]{IEEEtran}

\usepackage{cite}
\usepackage{amsmath}
\usepackage{booktabs}
\usepackage{array}
\usepackage{url}
\usepackage{hyperref}

\newcommand{\bench}{GPU-NFBench}

\begin{document}

\title{\bench: A Human-Labeled Benchmark and LLM-Assisted Triage System for GPU Numerical Failures}

\author{
\IEEEauthorblockN{Aryan Shah}
\IEEEauthorblockA{Independent Researcher\\
Texas, USA}
}

\maketitle

\begin{abstract}
GPU software stacks expose high-performance kernels through Python APIs,
array libraries, graph compilers, and domain-specific languages, but numerical
failures in these systems are difficult to diagnose from issue reports alone.
Reports about NaNs, overflow, dtype conversion, precision tolerance, and
incorrect outputs often mix user error, API limitations, compiler failures,
performance regressions, and true numerical defects. This paper presents
\bench, a human-labeled benchmark and triage framework for GPU numerical
failure reports. We construct an expanded 1,191-issue gold benchmark from
public GitHub issues across eight GPU-related repositories, using a seven-class
taxonomy that separates invalid-value failures, overflow/underflow,
precision/tolerance mismatches, dtype/casting failures, crash/compile
failures, performance-only reports, and non-numerical reports. We also preserve
the original 930-issue silver seed, the 191-issue adjudicated pilot benchmark,
and a full taxonomy-normalization log for the expanded labels. On the expanded
gold benchmark, weak candidate labels reach only 50.5\% accuracy, showing that
keyword retrieval is not sufficient ground truth. A deterministic full-coverage
voting ensemble reaches 73.8\% accuracy and 0.712 macro F1, while a fine-tuned
standalone FLAN-T5-base classifier reaches 80.5\% held-out test accuracy and
0.778 macro F1. Under leave-one-repository-out evaluation, the best
deterministic models reach 66.8\% accuracy, showing that project transfer
remains difficult. In a chronological split that trains on older reports and
tests on newer reports, the best deterministic model reaches 76.6\% accuracy
and 0.738 macro F1. Selective prediction further improves precision: the
deterministic ensemble reaches 84.2\% accuracy at 73.5\% coverage and 93.9\%
accuracy at 30.1\% coverage, while an LLM/deterministic agreement-abstention
mode reaches 89.3\% accuracy at 68.3\% coverage. A 119-row blind audit achieved
80.7\% Person A/B agreement and Cohen's $\kappa=0.721$; follow-on adjudication
applied 419 human label updates, changing 211 primary labels and producing the
canonical v2 benchmark used in the final experiments. We also provide a
250-row root-cause evidence-coded extension, including a 50-row
human-adjudicated subset, and an appendix of concrete boundary errors. These
results position \bench{} as a practical foundation for GPU-bug triage, dataset
construction, and future work on root-cause prediction for numerical failures
in ML systems.
\end{abstract}

\begin{IEEEkeywords}
GPU bugs, numerical bugs, benchmark construction, issue mining, LLM-assisted
triage, software reliability, machine learning systems
\end{IEEEkeywords}

\section{Introduction}

Modern machine-learning and scientific-computing workflows rely on GPU
software stacks that combine low-level kernels, compiler transformations, and
high-level Python interfaces. Frameworks such as Triton, PyTorch, JAX, CuPy,
Numba, RAPIDS, and TVM make GPU acceleration accessible, but they also create a
large surface for numerical failures. A small dtype mismatch, missing mask,
incorrect compiler lowering rule, or precision-tolerance assumption can
produce NaNs, Inf values, silent wrong answers, or failures that appear only on
specific hardware and software configurations.

Public issue trackers contain a large record of how these failures occur in
practice. However, issue reports are noisy. A search for ``NaN'' can return
true invalid-value bugs, missing-data examples, or unrelated benchmark tables.
A search for ``dtype'' can return casting failures, unsupported-feature
requests, performance complaints, or compiler crashes. These reports are
valuable, but only if the benchmark distinguishes candidate retrieval labels
from human gold labels.

This paper presents \bench, an expanded human-labeled benchmark for GPU
numerical failure reports. The benchmark is designed for practical triage:
given a public issue report, classify the dominant failure mode, determine
whether the issue is truly numerical, and optionally abstain when the evidence
is ambiguous. The artifact also preserves provenance for each stage: raw issue
collection, silver labels, pilot adjudication, expanded human labels, label
normalization, model training files, and evaluation reports.

The contributions are:
\begin{itemize}
    \item We construct an expanded 1,191-row human-labeled benchmark of GPU
    numerical-failure reports from public GitHub issues across eight
    GPU-related projects.
    \item We define a seven-class primary taxonomy and preserve a complete
    normalization log for 180 out-of-taxonomy labels in the expanded human
    annotation file.
    \item We evaluate weak labels, lexical baselines, nearest-neighbor
    retrieval, linear classifiers, a standalone fine-tuned FLAN-T5-base model,
    and a deterministic voting ensemble under stratified and
    leave-one-repository-out settings.
    \item We show that deterministic full-coverage triage reaches 73.8\%
    accuracy and 0.712 macro F1, while a standalone fine-tuned LLM reaches
    80.5\% held-out test accuracy and 0.778 macro F1.
    \item We evaluate selective triage: deterministic abstention reaches
    84.2\% accuracy at 73.5\% coverage, and an LLM/deterministic agreement
    mode reaches 89.3\% accuracy at 68.3\% coverage.
    \item We release a benchmark packet, data card, ablation tables, error
    analysis, linked-fix evidence, a completed blind audit, a 419-row
    adjudication update, a 50-row human-adjudicated root-cause subset, and a
    250-row evidence-coded root-cause extension.
\end{itemize}

\section{Background and Motivation}

GPU numerical bugs differ from ordinary application bugs because the symptom is
often detached from the underlying cause. A NaN in a model output can originate
from overflow in a fused operation, an unsafe dtype cast, an invalid memory
access, a compiler transformation, or a tolerance mismatch in the test suite.
Issue reports rarely state the root cause directly. Instead, they contain
partial evidence: a stack trace, a failing kernel, a device name, an expected
CPU result, a tolerance comparison, or a maintainer comment.

This makes the problem a good testbed for LLM-assisted software engineering.
Large language models can read issue text and infer likely failure categories,
but they also need calibrated taxonomies, strong negative classes, and
evaluation against human labels. In early experiments on the 191-issue pilot
benchmark, full-coverage models remained near 54\% accuracy, while selective
answering performed much better. The expanded benchmark tests whether adding a
larger human-labeled gold set makes full-coverage triage practical.

\section{Benchmark Construction}

\subsection{Seed Collection}

The project began with a 930-issue silver seed collected from six GPU-related
repositories: Triton, CuPy, PyTorch, JAX, Numba, and RAPIDS cuDF. Queries used
terms associated with numerical failure reports, including \texttt{nan},
\texttt{inf}, \texttt{overflow}, \texttt{dtype}, \texttt{precision},
\texttt{tolerance}, \texttt{incorrect}, and \texttt{wrong result}. Candidate
labels were assigned from query and text heuristics. This seed was useful for
coverage but was not treated as ground truth.

\subsection{Pilot Gold Benchmark}

From the seed, we constructed a 191-issue pilot gold benchmark with full public
issue/comment context. Two blind human annotation passes labeled primary
failure class, secondary causes, numerical-failure status, confidence, and
evidence snippets. Disagreements were adjudicated into final gold labels. The
pilot revealed two important facts: issue reports are genuinely ambiguous, and
weak labels are too noisy to serve as final supervision. Prior agreement
analysis on this pilot found modest overall primary-label agreement, but
substantially higher agreement on high-confidence cases. This motivated the
larger expanded benchmark and the selective-prediction setting.

\subsection{Expanded Human-Labeled Benchmark}

We expanded the benchmark with 1,000 additional completed human-labeled issue
records from an online collection queue. The expanded queue covers eight
repositories: JAX, PyTorch, RAPIDS cuML, Triton, RAPIDS cuDF, CuPy, Numba, and
Apache TVM. After merging the 1,000 expansion rows with the 191 pilot rows and
applying the v2 adjudication updates described below, the canonical expanded
gold benchmark contains 1,191 issues. Table~\ref{tab:labels} shows the
primary-label distribution.

\begin{table}[t]
\centering
\caption{Expanded gold benchmark label distribution.}
\label{tab:labels}
\begin{tabular}{lr}
\toprule
Primary label & Issues \\
\midrule
precision\_tolerance & 368 \\
dtype\_casting & 193 \\
crash\_compile & 192 \\
not\_numerical\_failure & 147 \\
nan\_inf & 143 \\
overflow\_underflow & 94 \\
performance\_only & 54 \\
\midrule
Total & 1191 \\
\bottomrule
\end{tabular}
\end{table}

\subsection{Taxonomy Normalization}

The expanded human annotation file included several labels outside the final
seven-class taxonomy, including API feature requests, compiler/code-generation
descriptions, generic needs-review labels, and performance-regression labels.
Rather than silently discard or recode them, we ran a deterministic
taxonomy-normalization pass and preserved the full change log. In total, 180
rows were mapped into the final taxonomy, and no invalid labels remained. For
example, performance-regression labels were mapped to \texttt{performance\_only};
compiler/code-generation labels were mapped to \texttt{crash\_compile},
\texttt{performance\_only}, or \texttt{not\_numerical\_failure} depending on
the recorded evidence; and API feature requests were mapped to
\texttt{not\_numerical\_failure}. The original labels remain recoverable from
the repair log, so this step is auditable.

\subsection{Expanded Agreement Audit and V2 Canonicalization}

The expanded 1,000-row addition was followed by a blind audit sampled from the
expansion set. The audit packet hides gold labels, candidate labels, model
predictions, and repair decisions. Two annotators independently labeled 119
rows using the final seven-class taxonomy. Person A and Person B agreed on
80.7\% of rows, with Cohen's $\kappa=0.721$, indicating substantial agreement
for a report-level taxonomy with overlapping symptom and root-cause cues.
A third adjudication pass then assigned final labels to all 119 audit rows.
The adjudicated audit differed from the existing expanded gold label in 52
rows, so we used it as a targeted revision signal rather than treating the
expanded labels as fixed. We then completed a second 300-row adjudication pass
focused on low-confidence and model-error examples. Together, the audit and
follow-on adjudication applied 419 human updates to the expanded benchmark and
changed 211 primary labels. The final experiments in this paper use this
canonical v2 benchmark. This revision step is important methodologically: it
turns disagreement analysis into benchmark repair rather than reporting label
noise as an unresolved limitation.

\section{Taxonomy}

Each issue receives exactly one primary label:
\begin{itemize}
    \item \textbf{nan\_inf}: invalid-value failures involving NaN, Inf, or
    non-finite outputs.
    \item \textbf{overflow\_underflow}: numerical range failures, including
    saturation, wraparound, or loss to zero.
    \item \textbf{precision\_tolerance}: wrong-answer or tolerance failures
    where approximate arithmetic, accumulation order, or precision is central.
    \item \textbf{dtype\_casting}: dtype promotion, conversion, unsupported
    dtype, or casting semantics failures.
    \item \textbf{crash\_compile}: failures dominated by crashes, compiler
    errors, illegal memory accesses, or build/runtime exceptions.
    \item \textbf{performance\_only}: reports about speed, memory use, or
    scheduling where no numerical correctness issue is present.
    \item \textbf{not\_numerical\_failure}: feature requests, API questions,
    documentation issues, and other non-numerical reports.
\end{itemize}

The taxonomy intentionally separates symptoms from suspected causes. For
example, an issue may be labeled \texttt{nan\_inf} as the primary symptom while
secondary evidence points to dtype conversion, masking, or compiler lowering.
This paper evaluates primary-label triage; secondary-cause prediction is left
as a follow-up task.

\section{Models and Evaluation}

\subsection{Evaluation Setup}

We evaluate models on the 1,191-row expanded gold benchmark. Metrics are
accuracy and macro F1 for the seven-class task. Macro F1 is important because
the dataset is imbalanced: precision/tolerance issues are the largest class,
while performance-only reports are the smallest. We also evaluate a binary
gate that predicts whether an issue is a numerical failure at all. The primary
headline setting is stratified five-fold evaluation. We also report
leave-one-repository-out evaluation, where each repository is held out as the
test project and all other repositories are used for training.

\subsection{Baselines}

The majority baseline predicts the largest class. The weak-label baseline uses
the original candidate label attached during retrieval. Text baselines include
BM25 nearest-neighbor retrieval, multinomial naive Bayes, TF-IDF logistic
regression, TF-IDF linear SVM, and bigram TF-IDF logistic regression. These
baselines are deliberately simple so that improvements can be attributed to
better supervision and calibrated decision rules rather than opaque modeling
choices.

\subsection{Standalone LLM}

We fine-tuned a standalone FLAN-T5-base sequence-to-sequence classifier on
benchmark-formatted issue text from the v2 canonical benchmark. The split
contains 945 training rows, 123 validation rows, and 123 held-out test rows. To
reduce class imbalance, the training set was balanced by oversampling minority
classes to 2,058 training examples. The selected v2 checkpoint starts from the
previous expanded-gold FLAN-T5-base checkpoint and is retrained for one epoch
on the v2 labels. It predicts one of the seven primary taxonomy labels directly
from issue text.

\subsection{Voting Ensemble and Abstention}

The best full-coverage deterministic system is a voting ensemble over the
strongest lexical and linear models. For selective triage, the ensemble can
abstain unless at least $k$ voters agree. This setting matches realistic
engineering use: a triage assistant should confidently label straightforward
issues and flag ambiguous reports for human review.

We also evaluate an LLM-assisted agreement mode. This mode answers only when
the deterministic ensemble and standalone LLM predict the same label; otherwise
it abstains. This is stricter than majority voting and is intended for
high-confidence triage queues where false labels are more costly than deferred
labels.

\subsection{Ablations and Error Analysis}

To test whether the ensemble depends too heavily on candidate labels, we report
ablations that remove candidate-label votes and retain only learned text
models. We also compute the most frequent gold/predicted error pairs and retain
representative issue examples for qualitative review.

\section{Results}

\subsection{Full-Coverage Classification}

Table~\ref{tab:fullcoverage} reports full-coverage performance. Weak candidate
labels reach only 50.5\% accuracy, confirming that keyword retrieval labels
are not reliable gold labels. Linear text models perform substantially better
with the expanded human-labeled set. The full-coverage voting ensemble reaches
73.8\% accuracy and 0.712 macro F1.

\begin{table}[t]
\centering
\caption{Expanded gold benchmark full-coverage results.}
\label{tab:fullcoverage}
\begin{tabular}{lcc}
\toprule
Model & Accuracy & Macro F1 \\
\midrule
Majority baseline & 0.309 & 0.067 \\
Candidate weak label & 0.505 & 0.426 \\
BM25 nearest neighbor & 0.571 & 0.534 \\
Naive Bayes & 0.610 & 0.501 \\
TF-IDF logistic regression & 0.715 & 0.688 \\
TF-IDF linear SVM & 0.732 & 0.703 \\
Bigram TF-IDF logistic & 0.724 & 0.694 \\
Voting ensemble & \textbf{0.738} & \textbf{0.712} \\
\bottomrule
\end{tabular}
\end{table}

\subsection{Cross-Repository Generalization}

Table~\ref{tab:loro} reports leave-one-repository-out results. This setting is
harder because model vocabularies, issue templates, and library-specific
concepts shift across projects. The best deterministic models in this setting
reach 66.8\% accuracy, with TF-IDF linear SVM reaching 0.648 macro F1 and the
voting ensemble reaching 0.647 macro F1. The result is substantially above the
weak-label baseline but lower than shuffled evaluation, indicating that
cross-project transfer remains a central challenge for GPU numerical-failure
triage.

\begin{table}[t]
\centering
\caption{Leave-one-repository-out evaluation on expanded gold.}
\label{tab:loro}
\begin{tabular}{lcc}
\toprule
Model & Accuracy & Macro F1 \\
\midrule
Candidate weak label & 0.505 & 0.426 \\
BM25 nearest neighbor & 0.493 & 0.478 \\
Naive Bayes & 0.550 & 0.475 \\
TF-IDF logistic regression & 0.653 & 0.630 \\
TF-IDF linear SVM & \textbf{0.668} & \textbf{0.648} \\
Bigram TF-IDF logistic & 0.663 & 0.633 \\
Voting ensemble & \textbf{0.668} & 0.647 \\
\bottomrule
\end{tabular}
\end{table}

Table~\ref{tab:repo_weakness} summarizes the weakest transfer repositories.
CuPy and Numba are the hardest held-out projects: CuPy mixes dtype/API
compatibility with numerical symptoms, while Numba contains CUDA, LLVM, build,
and runtime language that blurs \texttt{crash\_compile} and
\texttt{not\_numerical\_failure}. This supports cross-repository transfer as a
real benchmark challenge rather than a simple shuffled-text classification
task.

\begin{table}[t]
\centering
\caption{Weakest leave-one-repository-out transfer cases.}
\label{tab:repo_weakness}
\begin{tabular}{lccc}
\toprule
Held-out repo & Issues & Best acc. & Ensemble acc. \\
\midrule
numba/numba & 65 & 0.492 & 0.462 \\
cupy/cupy & 82 & 0.500 & 0.439 \\
rapidsai/cudf & 132 & 0.583 & 0.545 \\
\bottomrule
\end{tabular}
\end{table}

\subsection{Chronological Generalization}

We also evaluate a future-facing chronological split. Issues are sorted by
creation time; the older 80\% form the training set and the newer 20\% form the
test set. The split trains on 952 issues through December 17, 2025 and tests on
239 newer issues. Table~\ref{tab:time} shows that bigram TF-IDF logistic
regression reaches 76.6\% accuracy and 0.738 macro F1, while the voting
ensemble reaches 74.5\% accuracy and 0.722 macro F1. These results suggest
that the benchmark supports triage of later issue reports, but also that the
simpler linear model transfers better than the ensemble under temporal shift.

\begin{table}[t]
\centering
\caption{Chronological 80/20 evaluation on expanded gold.}
\label{tab:time}
\begin{tabular}{lcc}
\toprule
Model & Accuracy & Macro F1 \\
\midrule
Candidate weak label & 0.469 & 0.325 \\
Naive Bayes & 0.598 & 0.475 \\
TF-IDF logistic regression & 0.711 & 0.690 \\
TF-IDF linear SVM & 0.732 & 0.697 \\
Bigram TF-IDF logistic & \textbf{0.766} & \textbf{0.738} \\
Voting ensemble & 0.745 & 0.722 \\
\bottomrule
\end{tabular}
\end{table}

\subsection{Standalone LLM Performance}

Table~\ref{tab:llm} compares the standalone LLM with the best deterministic
ensemble. The v2 standalone FLAN-T5-base model reaches 80.5\% held-out test
accuracy and 0.778 macro F1, outperforming the deterministic ensemble on the
same held-out LLM test split. This changes the practical conclusion from the
earlier expanded-gold version: after adjudication repair, the standalone LLM is
not only a simpler deployable artifact, but also the strongest full-coverage
triage model in the final held-out evaluation. As an additional local
zero-shot comparison, we evaluated Ollama \texttt{llama3.2:3b} on the same 123
held-out rows. It reaches only 31.7\% accuracy and 0.322 macro F1, showing
that task-specific fine-tuning is important. A cloud-backed modern model was
not used for the final comparison because exporting held-out benchmark prompts
to an external service was blocked in the local environment.

\begin{table}[t]
\centering
\caption{Standalone LLM and deterministic comparison on v2 labels.}
\label{tab:llm}
\begin{tabular}{lcc}
\toprule
Model & Accuracy & Macro F1 \\
\midrule
Local Llama 3.2 3B zero-shot & 0.317 & 0.322 \\
TF-IDF linear SVM on LLM test split & 0.691 & 0.669 \\
Voting ensemble on LLM test split & 0.683 & 0.654 \\
FLAN-T5-base standalone LLM & \textbf{0.805} & \textbf{0.778} \\
Binary numerical gate, 5-fold & 0.892 & 0.796 \\
\bottomrule
\end{tabular}
\end{table}

\subsection{Selective Triage}

Table~\ref{tab:abstain} shows the accuracy/coverage tradeoff when the ensemble
abstains unless enough voters agree. With at least three agreeing voters, the
system labels 92.6\% of issues at 76.9\% accuracy. With at least four agreeing
voters, it labels 73.5\% of issues at 84.2\% accuracy. With unanimous or
near-unanimous agreement, it reaches 93.9\% accuracy on 30.1\% of issues. The
LLM/deterministic agreement mode provides a complementary high-confidence
setting: it answers 68.3\% of held-out LLM-test examples at 89.3\% accuracy
and 0.844 macro F1.

\begin{table}[t]
\centering
\caption{Selective prediction with abstention.}
\label{tab:abstain}
\begin{tabular}{lccc}
\toprule
Agreement rule & Coverage & Accuracy & Macro F1 \\
\midrule
Vote $\geq$ 2 & 0.999 & 0.739 & 0.713 \\
Vote $\geq$ 3 & 0.926 & 0.769 & 0.744 \\
Vote $\geq$ 4 & 0.735 & 0.842 & 0.803 \\
Vote $\geq$ 5 & 0.301 & 0.939 & 0.733 \\
LLM/ensemble agreement & 0.683 & 0.893 & 0.844 \\
\bottomrule
\end{tabular}
\end{table}

\subsection{Label-Signal Check}

The v2 results show that models are not merely replaying weak retrieval labels.
The candidate weak-label baseline reaches 50.5\% accuracy and 0.426 macro F1,
while the TF-IDF linear SVM reaches 73.2\% accuracy and 0.703 macro F1 and the
standalone LLM reaches 80.5\% held-out accuracy and 0.778 macro F1. The
improvement from weak labels to text-trained models indicates that the
human-labeled benchmark provides usable semantic supervision beyond query terms
or retrieval heuristics.

\subsection{Per-Class Behavior}

The deterministic ensemble is strongest on \texttt{precision\_tolerance},
\texttt{nan\_inf}, and \texttt{crash\_compile} issues, reaching F1 scores of
0.811, 0.769, and 0.740, respectively. It is weakest on
\texttt{performance\_only}, \texttt{not\_numerical\_failure}, and
\texttt{dtype\_casting}. These classes are boundary-heavy: performance reports
often include numerical benchmark terms, non-bug support questions can contain
stack traces or dtype language, and dtype changes can produce precision or
range symptoms.

\begin{table}[t]
\centering
\caption{Voting ensemble per-class performance.}
\label{tab:perclass}
\begin{tabular}{lccc}
\toprule
Class & Precision & Recall & F1 \\
\midrule
crash\_compile & 0.725 & 0.755 & 0.740 \\
dtype\_casting & 0.641 & 0.694 & 0.667 \\
nan\_inf & 0.855 & 0.699 & 0.769 \\
not\_numerical\_failure & 0.689 & 0.619 & 0.652 \\
overflow\_underflow & 0.813 & 0.649 & 0.722 \\
performance\_only & 0.744 & 0.537 & 0.624 \\
precision\_tolerance & 0.761 & 0.867 & 0.811 \\
\bottomrule
\end{tabular}
\end{table}

\subsection{Error Analysis}

The largest error families are semantically plausible boundary cases. The most
common confusion is \texttt{dtype\_casting} predicted as
\texttt{precision\_tolerance} (28 errors), followed by \texttt{nan\_inf}
predicted as \texttt{precision\_tolerance} (22 errors),
\texttt{precision\_tolerance} predicted as \texttt{dtype\_casting} (21 errors),
and \texttt{crash\_compile} predicted as \texttt{precision\_tolerance} (20
errors). These errors reflect the structure of GPU numerical reports: dtype
changes, compiler failures, overflow, and NaN symptoms are often reported
through failed numerical tolerance tests. The next largest errors involve
\texttt{not\_numerical\_failure} issues predicted as crash or dtype failures,
usually because issue text contains API type errors or installation stack
traces even when the report is not a numerical defect. This analysis supports
two follow-up tasks: multi-label secondary-cause prediction and explicit
separation of symptom labels from root-cause labels.

The artifact includes a 12-case error appendix with issue snippets and
boundary explanations. Examples include dtype reports predicted as
precision/tolerance because the failure appears as a numerical mismatch,
NaN/Inf reports embedded in broader tolerance failures, and non-bug support
requests predicted as crash or dtype failures because they contain stack-trace
or type-language cues.

\subsection{Qualitative Examples and Linked Fixes}

Representative examples clarify the label boundaries. A CuPy issue titled
``Wrong overflow with matmul of uint16 arrays'' is a direct
\texttt{overflow\_underflow} case. A Triton issue involving \texttt{remquo}
argument types is labeled \texttt{dtype\_casting} because the failure centers
on dtype signatures rather than numeric tolerance. A RAPIDS cuDF question about
a Python type error is labeled \texttt{not\_numerical\_failure}, even though it
contains type language, because the issue is an API question rather than a GPU
correctness failure.

The artifact also contains linked-fix evidence for future root-cause work. In
the pilot benchmark context, 105 issues contain at least one fix or root-cause
signal, including 39 explicit pull-request URLs. A follow-on public diff fetch
retrieved 488 changed-file rows from linked PRs. We additionally adjudicated a
50-row root-cause subset with five fix-oriented labels: compiler/backend/runtime
behavior, dtype or type semantics, numerical algorithm or tolerance,
performance/resource behavior, and non-bug feature/documentation/support cases.
To expand coverage without overstating provenance, we also provide a 250-row
evidence-coded root-cause extension. It contains the 50 human-adjudicated rows,
55 linked-fix evidence-coded rows, and 145 issue-evidence rows. Its largest
groups are numerical algorithm or tolerance cases (74), non-bug/support cases
(69), and dtype/type semantics cases (50). The paper treats this as
root-cause supervision and provenance for future work, not as 250
human-adjudicated root-cause labels.

\section{Discussion}

\subsection{Why the Expanded Benchmark Matters}

The expanded gold set changes the research conclusion. On the original
191-issue pilot, full-coverage classification remained near 54\%, and the most
useful setting was abstention. With 1,191 human-labeled examples, full-coverage
triage becomes viable: the deterministic ensemble reaches 73.8\% accuracy, the
standalone LLM reaches 80.5\% held-out accuracy, and chronological evaluation
reaches 76.6\% accuracy with the best deterministic model. This suggests that
the earlier ceiling was primarily a data and label-quality limitation, not
evidence that issue reports are inherently unclassifiable.

\subsection{Why Keep Both the LLM and Ensemble}

For a paper, the standalone LLM is now both the easiest deployable artifact and
the strongest full-coverage held-out model: it is one model that maps issue
text to a label. The deterministic ensemble remains useful as a transparent
baseline and confidence mechanism. Its abstention behavior is valuable because
ambiguous issue reports are common, and a practical triage assistant should
know when to defer. The strongest framing is therefore a two-level system: a
standalone LLM provides full-coverage semantic triage, while LLM/ensemble
agreement provides a high-confidence review queue.

\subsection{Practical Use Cases}

\bench{} can support three workflows. First, maintainers can prioritize issues
by separating likely numerical defects from performance or API reports. Second,
researchers can sample issue clusters for root-cause analysis, such as dtype
promotion failures or precision-tolerance mismatches. Third, future debugging
agents can use the benchmark as a supervised evaluation set before attempting
harder tasks such as patch localization or reproducer generation.

\section{Threats to Validity}

\textbf{Annotation quality.} The expanded benchmark relies on human labels and
a deterministic normalization pass. Although invalid labels were removed and
all repairs are logged, the full 1,191-row benchmark has not yet been
double-adjudicated. The 119-row audit shows strong A/B agreement, and the
follow-on v2 pass applied 419 adjudication updates, but the remaining
non-adjudicated rows can still contain boundary-label errors.

\textbf{Repository bias.} The benchmark is drawn from public GitHub issues in
popular GPU-related repositories. It may underrepresent private industrial
bugs, closed-source compiler failures, or hardware-specific issues that are not
reported publicly.

\textbf{Cross-project generalization.} We report leave-one-repository-out
results, but project sizes are uneven. Apache TVM has only 36 issues, while JAX
and PyTorch have nearly 300 each. Future versions should add more projects and
report time-based splits to test whether models generalize to new issue styles.

\textbf{Label granularity.} Some reports naturally span multiple categories.
For example, dtype conversion can lead to overflow, and compiler crashes can
hide underlying numerical defects. This paper uses one primary label for
evaluation, but multi-label secondary-cause modeling is an important extension.

\textbf{Issue-text limitation.} Many issue reports lack minimal reproducers,
environment metadata, or confirmed fixes. The benchmark classifies reports; it
does not prove the actual root cause unless linked fix evidence is available.

\section{Conclusion}

This paper presents \bench{}, a 1,191-row human-labeled benchmark for GPU
numerical failure reports. The expanded benchmark shows that reliable
full-coverage triage is possible when weak keyword labels are replaced with
human gold labels: a standalone FLAN-T5-base model reaches 80.5\% held-out
accuracy and 0.778 macro F1, a deterministic voting ensemble reaches 73.8\%
accuracy and 0.712 macro F1, chronological deterministic evaluation reaches
76.6\% accuracy and 0.738 macro F1, and LLM/deterministic agreement reaches
89.3\% accuracy at 68.3\% coverage. Leave-one-repository-out evaluation
reaches 66.8\% accuracy, showing that cross-project transfer remains difficult.
The remaining path to a stronger archival submission is clear: adjudicate more
of the expanded benchmark, add more projects, and extend report-level
classification into root-cause and fix-evidence tasks.

\section*{Artifact Availability}

The artifact contains the processed benchmark, expanded gold labels,
taxonomy-normalization report, model-training scripts, fine-tuning files,
evaluation tables, generated reports, qualitative examples, linked-fix evidence,
a data card, completed blind audit files, v2 adjudication files, LLM fine-tuning
files, the 50-row human-adjudicated root-cause subset, the 250-row
evidence-coded root-cause extension, cross-repository weakness analysis, local
LLM baseline outputs, a 12-case error appendix, a reproducibility checklist,
and an external-repository candidate pool for future benchmark expansion. The
data card documents intended use, non-use cases, label provenance, and known
limitations. The artifact is released at
\url{https://github.com/jubs-2431/gpu-nfbench/releases/tag/v1.0-conference}.
The archived Zenodo record is available at
\url{https://zenodo.org/records/20242157} with DOI
\url{https://doi.org/10.5281/zenodo.20242157}.

\begin{thebibliography}{9}

\bibitem{triton2019}
P. Tillet, H. Kung, and D. Cox, ``Triton: An intermediate language and compiler
for tiled neural network computations,'' in \emph{Proceedings of the 3rd ACM
SIGPLAN International Workshop on Machine Learning and Programming Languages},
2019.

\bibitem{cupy2017}
R. Okuta, Y. Unno, D. Nishino, S. Hido, and C. Loomis, ``CuPy: A NumPy-compatible
library for NVIDIA GPU calculations,'' in \emph{Proceedings of Workshop on
Machine Learning Systems at NeurIPS}, 2017.

\bibitem{pytorch2019}
A. Paszke et al., ``PyTorch: An imperative style, high-performance deep learning
library,'' in \emph{Advances in Neural Information Processing Systems}, 2019.

\bibitem{jax2018}
J. Bradbury et al., ``JAX: Composable transformations of Python+NumPy programs,''
2018. [Online]. Available: \url{https://github.com/jax-ml/jax}

\bibitem{numba2015}
S. K. Lam, A. Pitrou, and S. Seibert, ``Numba: A LLVM-based Python JIT
compiler,'' in \emph{Proceedings of the Second Workshop on the LLVM Compiler
Infrastructure in HPC}, 2015.

\bibitem{rapids2024}
RAPIDS Development Team, ``RAPIDS: Collection of GPU accelerated data science
libraries,'' 2024. [Online]. Available: \url{https://rapids.ai/}

\bibitem{tvm2018}
T. Chen et al., ``TVM: An automated end-to-end optimizing compiler for deep
learning,'' in \emph{13th USENIX Symposium on Operating Systems Design and
Implementation}, 2018.

\bibitem{t5}
C. Raffel et al., ``Exploring the limits of transfer learning with a unified
text-to-text transformer,'' \emph{Journal of Machine Learning Research}, vol.
21, no. 140, pp. 1--67, 2020.

\bibitem{flan}
H. W. Chung et al., ``Scaling instruction-finetuned language models,'' \emph{J.
Mach. Learn. Res.}, vol. 25, no. 70, pp. 1--53, 2024.

\end{thebibliography}

\end{document}
